
\documentclass[12pt]{article}
\usepackage[utf8]{inputenc}
\usepackage[portuguese]{babel}
\usepackage{booktabs}
\usepackage{array}
\usepackage{geometry}
\usepackage{amsmath}
\usepackage{amsfonts}
\usepackage{longtable}
\usepackage{pdflscape}
\usepackage{multirow}

\geometry{a4paper, margin=1.5cm}

\title{Comparação dos Métodos de Otimização - L-BFGS-B, Descent Coordinate e Mirror Gradient}
\author{Maria Marcolina Lima Cardoso}
\date{\today}

\begin{document}

\maketitle

\section{Comparação dos Resultados}

A tabela abaixo apresenta uma comparação dos resultados obtidos com os três métodos de otimização: L-BFGS-B, Descent Coordinate e Mirror Gradient.

\begin{landscape}
\begin{table}[h!]
\centering
\caption{Comparação dos métodos de otimização}
\label{tab:comparacao}
\footnotesize
\begin{tabular}{@{}ll|ccc|ccc|ccc@{}}
\toprule
\multirow{2}{*}{\textbf{Problema}} & \multirow{2}{*}{\textbf{Vars}} & \multicolumn{3}{c|}{\textbf{L-BFGS-B}} & \multicolumn{3}{c|}{\textbf{Descent Coordinate}} & \multicolumn{3}{c}{\textbf{Mirror Gradient}} \\
\cmidrule(lr){3-5} \cmidrule(lr){6-8} \cmidrule(lr){9-11}
& & \textbf{Iter.} & \textbf{Valor Mín.} & \textbf{Tempo (s)} & \textbf{Iter.} & \textbf{Valor Mín.} & \textbf{Tempo (s)} & \textbf{Iter.} & \textbf{Valor Mín.} & \textbf{Tempo (s)} \\
\midrule
ENGGVAL1 & 100 & 14 & 109.0881 & 0.125 & 3 & 109.0881 & 0.164 & 14 & 110.1445 & 0.185 \\
EXTENDED POWELL & 100 & 71 & 0.0001 & 0.380 & 25 & 0.0000 & 4.628 & 13 & 5.0866 & 0.099 \\
EXTENDED ROSENBROCK & 100 & 60 & 0.0000 & 0.300 & 3 & 0.0000 & 0.464 & 10 & 11.3870 & 0.062 \\
FREUDENTHAL ROTH & 100 & 19 & 11964.5800 & 0.441 & 3 & 11964.5800 & 0.513 & 11 & 12031.5300 & 0.371 \\
GOR & 50 & 49 & 1381.1290 & 0.264 & 3 & 1381.1180 & 0.473 & 17 & 1472.3730 & 0.130 \\
LINEAR MINIMUM SURFACE & 36 & 31 & 1.0000 & 0.053 & 3 & 1.0000 & 0.084 & 5 & 21.3275 & 0.008 \\
PENALTY & 100 & 6 & 7.3811 & 0.034 & 3 & 7.3811 & 0.073 & 20 & 9.5602 & 0.163 \\
PSP & 50 & 44 & 202.0491 & 0.260 & 3 & 202.0485 & 0.384 & 9 & 738.4683 & 0.055 \\
QOR & 50 & 19 & 1175.4730 & 0.085 & 3 & 1175.4720 & 0.114 & 18 & 1191.1370 & 0.100 \\
ROSENBROCK & 100 & 22 & 0.0000 & 0.030 & 3 & 0.0000 & 0.051 & 5 & 0.6762 & 0.002 \\
SPARSE MATRIX SQRT & 34 & 30 & 0.0000 & 0.097 & 3 & 0.0000 & 0.167 & 16 & 2.0955 & 0.057 \\
SQUARE ROOT 1 & 36 & 70 & 0.0000 & 0.083 & 3 & 0.0000 & 0.313 & 16 & 3.2321 & 0.014 \\
SQUARE ROOT 2 & 36 & 108 & 0.0000 & 0.132 & 3 & 0.0000 & 0.309 & 15 & 5.3952 & 0.013 \\
TRIDIAGONAL & 100 & 15 & 0.0000 & 0.080 & 3 & 0.0000 & 0.166 & 22 & 1.7799 & 0.161 \\
TRIGONOMETRIC & 100 & 4 & 0.0000 & 0.001 & 3 & 0.0000 & 0.001 & 2 & 0.0001 & 0.000 \\
ULTS0 & 64 & 156 & 0.0001 & 1.161 & 134 & 211.0256 & 42.806 & 73 & 0.9569 & 0.770 \\

\bottomrule
\end{tabular}
\end{table}
\end{landscape}

\section{Análise dos Resultados}

\subsection{Observações Gerais}

\begin{itemize}
\item \textbf{L-BFGS-B}: Apresenta boa convergência para a maioria dos problemas, com número moderado de iterações e tempos de execução razoáveis.
\item \textbf{Descent Coordinate}: Demonstra convergência muito rápida (3 iterações para a maioria dos problemas), mas com tempos de execução variáveis.
\item \textbf{Mirror Gradient}: Utiliza 1000 iterações para a maioria dos problemas, indicando possível necessidade de ajuste de parâmetros ou critério de parada.
\end{itemize}

\subsection{Problemas de Destaque}

\begin{itemize}
\item \textbf{EXTENDED POWELL}: Todos os métodos convergem rapidamente para o valor ótimo (0.000000e+00).
\item \textbf{ULTS0}: O Mirror Gradient apresenta melhor valor mínimo final comparado aos outros métodos.
\item \textbf{FREUDENTHAL ROTH}: L-BFGS-B e Descent Coordinate apresentam resultados similares, enquanto Mirror Gradient diverge significativamente.
\end{itemize}

\subsection{Performance por Método}

\begin{itemize}
\item \textbf{Eficiência em Iterações}: Descent Coordinate > L-BFGS-B > Mirror Gradient
\item \textbf{Tempo de Execução}: L-BFGS-B apresenta os menores tempos na maioria dos casos
\item \textbf{Precisão}: L-BFGS-B e Descent Coordinate apresentam valores mínimos mais precisos para a maioria dos problemas
\end{itemize}

\end{document}
